\section{A Constitution for a More Capable Civilization}

The four principles concern boundaries, correction, division of work, and continuing revision. At a larger scale, they ask where a civilization will direct the capacity it has gained.

\subsection{Keep a Frontier That Answers Back}

Some problems compel sustained contact with material limits: deep-space missions, ocean engineering, ecological restoration, fundamental experiments, infrastructure renewal, and the formation of capable people.

They demand resources, time, and coordination. A useful shorthand for projects with high requirements on all three coordinates is
\begin{equation}
 \mathcal F=
 \{(R,T,G):R\geq R_\ast,\ T\geq T_\ast,\ G\geq G_\ast\}.
\end{equation}
This defines the intersection chosen by those thresholds. A broader class of difficult projects may be high in only one or two coordinates.

What makes such projects worth doing is their benefit, knowledge, or protection of future possibilities. Their resistance ensures that a convincing description cannot substitute for the required result.

A spacecraft needs an actual trajectory and reliable systems. An ocean instrument needs to survive pressure and corrosion. An ecosystem intervention needs observations through changing seasons and responses. A new scientific claim may need an instrument that does not yet exist.

Feedback is often costly and delayed in these settings. That can make evidence sparse rather than unusually dense. A well-designed program buys discriminating tests, preserves failure records, and stages commitments so uncertainty can be reduced before the largest exposures.

AI can improve design, analysis, and coordination across these frontiers. The ambition is to obtain more knowledge and service from each costly contact. Calling a problem “physical” does not make its current method optimal.

Long projects need a public account of opportunity cost. Resources directed toward a mission are unavailable elsewhere. Current decision-makers must explain what future benefit they expect, who bears the interim burden, and how later participants can revise the commitment.

Some benefits are shared or delayed enough that individual investors will not fund the socially useful amount. Public institutions, cooperative arrangements, and long-lived funds can carry the account across those periods.

\subsection{Life, Mind, and Trust as Areas of Growing Demand}

Three domains concentrate needs that cheap generation alone does not meet: maintaining bodies, supporting mental and relational life, and establishing dependable judgments.

Biotechnology can change what is treatable. Programmable genome editing supplied a major mechanism \cite{jinek2012programmable}. Clinical delivery, safety, access, and long-term evidence determine which benefits reach people. The 2026 Casgevy record is one current example of progress joined to continued obligations \cite{fda2026casgevy}.

Mental health requires similarly careful boundaries. WHO's updated 2025 report describes more than a billion people living with mental disorders and substantial gaps in services \cite{who2025mental}. A conversational interface may provide useful support. Clinical effectiveness, safety, escalation, privacy, and access need evidence appropriate to the service advertised.

Trust becomes a technical and institutional product when convincing signals are cheap. Provenance, secure identity, independent checking, and enforceable responsibility can help reconnect claims to their origins and consequences. C2PA illustrates a defined technical contribution, with its validation scope made explicit \cite{c2pa2026spec}.

These needs can expand industries. They can also remain poorly paid or unavailable to people who need them most. A demand forecast does not supply a distribution system.

Baumol's mechanism clarifies one possible relative-price pressure \cite{baumol1967macro}. If the progressive sector's productivity grows at constant rate $r$, the slower sector's does not, and wages move together under the model, then a relative unit-cost index can evolve as
\begin{equation}
 \frac{C_{\mathrm{slow}}(t)}{C_{\mathrm{fast}}(t)}
 =
 \frac{C_{\mathrm{slow}}(0)}{C_{\mathrm{fast}}(0)}
 e^{rt}.
\end{equation}
The sector outputs must be consistently indexed; they are not literally the same product. Different productivity, wage, capital, and service-quality paths change the result.

Rising relative costs can generate investment, shortages, unpaid family burdens, or exclusion. Public payment, insurance, labor institutions, and service design influence which outcome occurs.

\subsection{The Gaia Hypothesis as a Design Question}

A more speculative possibility follows from coupling industrial and biological systems: could a civilization develop stronger capacities to maintain, repair, and reproduce its material basis?

Von Neumann's self-reproducing automata separated a description from machinery able to use and copy it within a formal construction \cite{vonneumann1966theory}. Biological replication adds actual materials, energy, error correction, and selection. Engineering reproduction depends on supply chains whose independence has to be demonstrated.

Self-maintenance, reproduction, and heritable variation are useful properties to inspect. Their presence at some scale does not establish that the entire planet is one organism or that an industrial network can reproduce autonomously from raw resources.

Copy fidelity gives one measurable constraint. With independent per-site error probability $\nu$ over length $\lambda$,
\begin{equation}
 Q=(1-\nu)^\lambda\approx e^{-\nu\lambda}.
\end{equation}
The approximation requires small per-site error, with accumulated higher-order terms negligible at the desired accuracy. Selection, repair, redundancy, and the effects of particular mutations determine what information persists \cite{kunkel2000dna,gilbert1986}.

Circular material use provides another measure. For a defined material, boundary, and period,
\begin{equation}
 \kappa_{\mathrm{mat}}=
 \frac{\dot m_{\mathrm{returned\ to\ usable\ supply}}}
      {\dot m_{\mathrm{total\ usable\ input}}}.
\end{equation}
The numerator must be part of the denominator for the fraction to lie between zero and one. Stocks, delays, material quality, and repeated counting need explicit treatment.

A high recycling share can reduce demand for new extraction. It does not establish energetic self-sufficiency, complete reproduction, or ecological stability. Recycling uses energy, loses or disperses some material, and may change quality. An open system still needs energy throughput and entropy export.

“Gaia” is useful here as a question about coupled self-maintenance. It provides no guarantee that the maintained state would favor humans. A stable planetary process can be hostile to a particular species.

The engineering frontier is consequently specific: improve repair, circularity, resource resilience, and compatible ecological relations while preserving control over the consequences. A self-maintaining subsystem may be easier to sustain and harder to stop once dispersed. Both belong in its evaluation.

\subsection{Four Commitments}

A practical constitution for this setting can begin with four commitments.

\textbf{Keep correction possible.} Evidence, appeal, revision, and interruption must remain available as capability grows.

\textbf{Give affected people a voice in the cost.} Exposure should bring information, representation, and remedies, especially where burdens can be exported.

\textbf{Preserve necessary functions while improving their means.} Learning, verification, consent, recovery, and responsibility should survive the removal of waste.

\textbf{Evaluate the whole consequence over the relevant time.} A fast interface should not claim a gain whose bill sits in an unobserved body, network, or future.

Silicon--carbon symbiosis names one possible arrangement of these commitments. It is a shorthand for complementary capabilities and stakes, not a permanent division based on the periodic table. Digital systems remain material; living people also learn, compute, copy, and delegate.

The desired arrangement lets powerful tools expand human and ecological possibilities while keeping the consequences governable. It has to be built through standards, incentives, evidence, rights, and resources.

The final chapter tests the framework against two extremes: systems that reproduce their own physical basis, and systems that seem to replace physical life with representation. Both return to the same question: what must keep running for the possibility to remain real?
